% Qualité de la recherche : les deux métriques sont bornées entre 0 et 1 et
% partagent donc un axe unique. Aucun axe secondaire.
%
% Barres horizontales, comme les autres graphiques du chapitre : l'étiquette
% de valeur se place au bout de chaque barre, ce qui rend toute collision
% impossible — en barres verticales, « 0,969 » et « 0,891 » se touchaient.
% Recall@5 et MRR par méthode de recherche
% Généré par rapport/build_data.py — NE PAS ÉDITER.
\newcommand{\RechercheRecall}{(0.849,0) (0.660,1) (0.868,2)}
\newcommand{\RechercheMrr}{(0.778,0) (0.535,1) (0.775,2)}

\begin{tikzpicture}
\begin{axis}[
  socle,
  xbar, bar width=8pt,
  height=5.2cm, width=0.72\linewidth,
  xmin=0, xmax=1.18,
  xtick={0,0.25,0.5,0.75,1},
  xmajorgrids, ymajorgrids=false,
  axis x line*=bottom, axis y line*=left,
  y axis line style={draw=none}, ytick style={draw=none},
  ytick={0,1,2}, y dir=reverse,
  yticklabels={{Dense seul},{BM25 seul},{Hybride (RRF)}},
  yticklabel style={font=\sffamily\small, text=ink, align=right},
  enlarge y limits=0.28,
  xlabel={valeur de la métrique},
  nodes near coords,
  % « fixed zerofill » : sans lui, 0,660 s'affiche « 0,66 » alors que le
  % tableau voisin écrit trois décimales.
  nodes near coords style={font=\sffamily\scriptsize, text=ink,
                           anchor=west, xshift=2.5pt,
                           /pgf/number format/fixed,
                           /pgf/number format/fixed zerofill,
                           /pgf/number format/precision=3},
  legend style={at={(0.5,1.03)}, anchor=south},
  legend image post style={line width=2.5pt},
]
% L'axe est inversé (y dir=reverse) : un décalage positif place la barre plus
% haut. Recall@5 vient donc en premier, dans l'ordre de la légende.
\addplot[draw=inptblue,  fill=inptblue,  fill opacity=0.9,  bar shift=5pt]
  coordinates {\RechercheRecall};
\addplot[draw=amberdark, fill=amberdark, fill opacity=0.9,  bar shift=-5pt]
  coordinates {\RechercheMrr};
\legend{Recall@5, MRR}
\end{axis}
\end{tikzpicture}
